\documentclass[10pt,a4paper]{amsart}
\usepackage[margin=19mm]{geometry}
\usepackage{amsmath,amssymb,mathtools}
\usepackage[scr=rsfs,cal=euler]{mathalfa}
\usepackage{xfrac}
\usepackage[unicode,hidelinks,bookmarks=false]{hyperref}
\newcommand{\ie}{i.e.\ }
\newcommand{\cf}{cf.\ }
\newcommand{\resp}{respectively\ }
\newcommand{\Subsection}{\subsection}
\providecommand{\nobreakdash}{\nobreak\hbox{-}}
\setlength{\emergencystretch}{2em}
\multlinegap0pt
\usepackage{bm}
\usepackage{bbm}
\usepackage{dsfont}
\usepackage{xspace}
\usepackage{cancel}
\usepackage{mathtools}
\usepackage{amsthm}
\makeatletter
\@ifundefined{theorem}{\newtheorem{theorem}{Theorem}}{}
\@ifundefined{proposition}{\newtheorem{proposition}[theorem]{Proposition}}{}
\makeatother
\def\ge{\geqslant}
\def\le{\leqslant}
\newcommand{\sgn}{\mathrm{sgn}}
\renewcommand{\det}{\mathrm{det}}
\title[On AAA approximants on an interval to Nevanlinna-Pick functi]{On AAA approximants on an interval to Nevanlinna-Pick functions and P\'olya Frequency Series}
\author{Sergey A. Denisov, Maxim L. Yattselev}
\date{Source: arXiv:2608.24611v1 (CC BY 4.0)}
\begin{document}
\maketitle

\noindent\emph{Source and licence.} On AAA approximants on an interval to Nevanlinna-Pick functions and P\'olya Frequency Series. Version arXiv:2608.24611v1 (CC BY 4.0), distributed under the Creative Commons Attribution 4.0 International licence, as recorded in the arXiv licence metadata for this exact version and shown on the abstract page https://arxiv.org/abs/2608.24611v1. Full rights notice: licenses/arxiv-2608.24611-CC-BY-4.0.txt.\par\smallskip

\noindent\textit{Editorial scope.} Counted unit: the Proposition whose source label is \texttt{prop:esr} in the article (source0.tex lines 526--529, source label \texttt{prop:esr}), delivered together with the article's own complete original proof of it (source0.tex lines 530--606, one \texttt{proof} environment of 77 lines). No mathematics is written, repaired or re-derived: statement and proof are the article's own text line for line. Required context retained uncounted: PFS-Taylor (lines 197--200); PFS (lines 202--205); der K (lines 505--508). Every definition, display and label of those lines is retained unchanged. Omitted from this package and not redistributed: 1--196, 201--201, 206--504, 509--525, 607--694. Nothing inside a retained excerpt is omitted or altered: every retained body is delivered exactly as the article's lines are written, so no omission receipt of any kind accompanies this package. Citations: the retained excerpts carry no citation command at all, so no bibliography entry of the article is transcribed and no reference is redistributed. Reference adaptation: none was needed and none was made; every reference inside the delivered unit resolves inside it, so no unresolved reference or citation is produced. Printed slips kept literal: no printed word, symbol, hypothesis or number of the counted statement or of its proof was altered. Layout and class: the article's own class is \texttt{amsart} with packages including pdfsync, xcolor, hyperref, epsfig, graphicx, subfigure, amsmath, amssymb, newtxtext, newtxmath, mathabx, mathtools; the wrapper is the lane's pinned \texttt{amsart} editorial wrapper at 10pt on A4, which restates the theorem-like environments the retained text needs and adds the article's own macro definitions that the retained text needs (\texttt{\textbackslash det}, \texttt{\textbackslash ge}, \texttt{\textbackslash le}, \texttt{\textbackslash sgn}), with extra wrapper packages (bm, bbm, dsfont, xspace, cancel, mathtools). The delivered numbering is the unit's own. Assets: no retained span contains a figure environment, a graphics-inclusion command, a PDF-page inclusion, a table or a TikZ picture; no image file is copied into this package, the package declares no asset dependency and no image file is redistributed with it. Licence: version arXiv:2608.24611v1 of this article is distributed under the Creative Commons Attribution 4.0 International licence, as recorded in the arXiv licence metadata for this exact version and shown on the abstract page https://arxiv.org/abs/2608.24611v1. A full rights notice is delivered at licenses/arxiv-2608.24611-CC-BY-4.0.txt. Characters: every retained excerpt is pure ASCII, so the delivered engine, XeLaTeX with its default Latin Modern text font, reports no missing character.
\begin{align}
\label{PFS-Taylor}
F(z) = a_0 + a_1z + \cdots + a_nz^n + \cdots,
\end{align}

\begin{align}
\label{PFS}
F(z) = cz^\lambda e^{\gamma z}\prod_{i=1}^\infty \frac{1+\alpha_iz}{1-\beta_i z},
\end{align}

\begin{align}
\label{der K}
K_f^{(i,j)}(x,x) = \frac{i!j!}{(i+j+1)!} f^{(i+j+1)}(x) = i!j! a_{i+j+1}(x),
\end{align}

\begin{proposition}
\label{prop:esr}
Let \( F(z) \) be a non-rational P\'olya frequency series \eqref{PFS}. Further, let \( X \) be an open interval between the largest zero of \( F(z) \) and its smallest pole (\( X \) is unbounded on the left if \( F(z) \) has no zeros and is unbounded on the right if \( F(z) \) has no poles). Then, the Loewner kernel \( K_F(x,y) \) is extended sign-regular on \( X\times X \) with the sign pattern \( +--++--++--\cdots\). 
\end{proposition}

\begin{proof}
Let \( r_{\mathrm z} \) be the largest zero of \( F(z) \) (necessarily \( r_{\mathrm z}\leq0 \)) or an arbitrary negative number if \( F(z) \) has no zeros. As mentioned before the proposition, \( F(z+r_{\mathrm z})\) is also a P\'olya frequency series. Hence, it is enough to prove the desired claim for \( X=(0,r_{\mathrm p})\), where \( r_{\mathrm p} \) is the smallest pole of \( F(z) \) and necessarily is also the radius of convergence of the series~\eqref{PFS-Taylor} (\( r_{\mathrm p}=\infty \) if \( F(z)\) has no poles). Fix a natural number \( l\geq 1 \). By definition, 
\begin{align}
\label{W det}
W_l(x,y) = \det\left( K^{(i-1,j-1)}(x,y)\right)_{i,j=1}^l = \sum_\sigma (-1)^\sigma \prod_{i=1}^l K^{(i-1,\sigma(i)-1)}(x,y),
\end{align}
where the sum is taken over all the permutations \( \sigma \) of \( \{1,2,\ldots,l\} \) and \( (-1)^\sigma \) denotes the sign of the permutation. It follows from the Taylor theorem for analytic functions in two variables that 
\begin{align}
\label{W Taylor}
W_l(x,y) = \sum_{n,m=0}^\infty w_{l;n,m}x^ny^m, \quad w_{l;n,m} = \frac{W_l^{(n,m)}(0,0)}{n!m!},
\end{align}
where the series is convergent for all \( |x|,|y|<r_{\mathrm p} \). To prove the proposition, it is sufficient to show that the above Taylor coefficients \( w_{l;n,m} \) all have the required sign pattern: $\sgn \,w_{l;n,m}=(-1)^{l(l-1)/2}$. We apply the generalized Leibniz rule to \eqref{W det} to conclude  that \( W_l^{(n,m)}(x,y) \) is equal to
\begin{align}
\sum_{|\boldsymbol q|=m}\sum_{|\boldsymbol p|=n}\frac{m!}{q_1!\cdots q_l!}\frac{n!}{p_1!\cdots p_l!}\sum_\sigma (-1)^\sigma \prod_{i=1}^l K^{(p_i+i-1, q_{\sigma(i)} + \sigma(i) -1)}(x,y),
\end{align}
where \( \boldsymbol p=(p_1,\ldots,p_l) \), \( |\boldsymbol p|=p_1 + \cdots + p_l\), and \( p_i\geq 0 \) (the same notation holds for \( \boldsymbol q\)). Let, as before, \( a_n \) be the Taylor coefficients of \( F(z) \) at the origin. By evaluating the above expression at \( (0,0) \), dividing it by \( n!m! \), and using \eqref{der K}, we obtain
\begin{align}
w_{l;n,m} &= \sum_{|\boldsymbol q|=m}\sum_{|\boldsymbol p|=n} \left(\prod_{s=1}^l \frac{(p_s+s-1)!}{p_s!}\frac{(q_s+s-1)!}{q_s!}\right) \det\big(a_{p_i+i+q_j+j-1}\big)_{i,j=1}^l \\
\label{wnm}
& = \sum_{|\boldsymbol q|=m}\sum_{|\boldsymbol p|=n} \left(\prod_{s=1}^l \frac{\Gamma(p_s+s)}{\Gamma(p_s+1)}\frac{\Gamma(q_s+s)}{\Gamma(q_s+1)}\right) \det\big(a_{p_i+i+q_j+j-1}\big)_{i,j=1}^l,
\end{align}
where convenience of the second representation will become apparent momentarily. To proceed,  notice that, up to a sign, several of the summands above correspond to the same Hankel determinant; that is, their corresponding matrices differ only by row and/or column permutations. Hence, our next goal is to bring all determinants that are identical up to sign into a standard (Hankel) form. Suppose first that \( p_i+i=p_k+k \) for some \( i \) and \( k, i\neq k \), then the corresponding determinant vanishes since the underlying matrix has two identical rows. Otherwise, let \( \boldsymbol p \) be such that all \( p_i+i \) are all distinct. Then, there exists a permutation \( \sigma_{\boldsymbol p} \) such that
\[
p_{\sigma_{\boldsymbol p}(1)}+\sigma_{\boldsymbol p}(1)<p_{\sigma_{\boldsymbol p}(2)}+\sigma_{\boldsymbol p}(2)< \cdots < p_{\sigma_{\boldsymbol p}(l)}+\sigma_{\boldsymbol p}(l).
\]
Define \( u_i = p_{\sigma_{\boldsymbol p}(i)}+\sigma_{\boldsymbol p}(i)-i \). It holds that \( u_1 = p_{\sigma_{\boldsymbol p}(1)}+\sigma_{\boldsymbol p}(1)-1\geq p_{\sigma_{\boldsymbol p}(1)} \geq 0 \) and
\[
u_{i+1} = p_{\sigma_{\boldsymbol p}(i+1)}+\sigma_{\boldsymbol p}(i+1)-i - 1 > p_{\sigma_{\boldsymbol p}(i)}+\sigma_{\boldsymbol p}(i)-i -1 = u_i - 1.
\]
Hence, \( u_{i+1}\geq u_i\geq \ldots \ge u_1\geq 0 \). Observe also that
$
|\boldsymbol u|=|\boldsymbol p|=n\,.
$
Denote by $\boldsymbol \Upsilon$ the collection of all vectors $\boldsymbol u=(u_1,\ldots,u_l)$ with integer coordinates that satisfy the properties that $|\boldsymbol u|=n$ and {$0\le u_1\le \ldots \le u_{l-1}\le u_{l}$.} Our argument above shows that every $\boldsymbol p$ corresponds to some element in $\boldsymbol \Upsilon$. If $\boldsymbol u$ is fixed and $\boldsymbol p$ gives rise to such $\boldsymbol u$, we will write 
\( \boldsymbol p\sim \boldsymbol u \). Conversely, suppose \( \boldsymbol u \in \boldsymbol \Upsilon\). Given a permutation \( \sigma \) of \( \{1,\ldots,l\} \), set
\[
p_{\sigma(i)}^\sigma := u_i+i-\sigma(i), \quad i\in\{1,\ldots,l\}.
\]
It is still true that \( p_1^\sigma+\cdots +p_l^\sigma = n \) and that \( p_i^\sigma+i>0 \), but, in general, it is not true that \( p_i^\sigma\geq0 \). However, if \( p_i^\sigma<0 \), then \( \Gamma^{-1}(p_i^\sigma+1)=0 \). Denote \( D(\boldsymbol p,\boldsymbol q) := \det\big(a_{p_i+i+q_j+j-1}\big)_{i,j=1}^l \). Therefore, for fixed $\boldsymbol u$ and $\boldsymbol q$, we get 
\begin{align}
\sum_{\boldsymbol p\sim \boldsymbol u} \left(\prod_{s=1}^l\frac{\Gamma(p_s+s)}{\Gamma(p_s+1)}\right)D(\boldsymbol p,\boldsymbol q)& = \left(\sum_\sigma (-1)^\sigma \prod_{s=1}^l\frac{\Gamma(p_s^\sigma+s)}{\Gamma(p_s^\sigma+1)}\right) D(\boldsymbol u,\boldsymbol q) \\
& = \left(\sum_\sigma (-1)^\sigma \prod_{i=1}^l \frac{\Gamma(u_i+i)}{\Gamma(u_i+i-\sigma(i)+1)}\right) D(\boldsymbol u,\boldsymbol q).
\end{align}
Observe that
\begin{eqnarray*}
W(\boldsymbol u) := 
\sum_\sigma (-1)^\sigma \prod_{i=1}^l \frac{\Gamma(u_i+i)}{\Gamma(u_i+i-\sigma(i)+1)} =
\det\left( \frac{\Gamma(u_i+i)}{\Gamma(u_i+i-k+1)} \right)_{i,k=1}^l.
\end{eqnarray*}
Writing this explicitly in the matrix form gives
\[
W(\boldsymbol u) = \det\begin{pmatrix}
1 & u_1 & \cdots & u_1(u_1-1)\cdots(u_1-l+2) \\
1 & u_2+1 & \cdots & (u_2+1)(u_2)\cdots(u_2-l+3) \\
\vdots & \vdots & \ddots & \vdots \\
1 & u_l+l-1 & \cdots & (u_l+l-1)(u_l+l)\cdots(u_l+1)
\end{pmatrix}.
\]
Upon close examination it becomes clear that the above determinant can be rewritten as the Vandermonde determinant in \( u_1,u_2+1,\ldots,u_l+l-1 \). Hence, since the entries in $\boldsymbol u$ are increasing, we have that
\begin{align}
\label{sign Wu}
W(\boldsymbol u) = \prod_{1\leq i<j\leq l} (u_j-u_i+j-i)>0.
\end{align}
Exactly the same considerations apply to the sum over \( \boldsymbol q \) in \eqref{wnm}. Thus, we get that
\[
w_{l;n,m} = \sum_{\boldsymbol v\in \boldsymbol\Upsilon}\sum_{\boldsymbol u\in \boldsymbol\Upsilon} W(\boldsymbol v)W(\boldsymbol u) D(\boldsymbol u,\boldsymbol v).
\]
Since \( F(z) \) is a non-rational P\'olya frequency series, it then holds that
\begin{align}
\label{sign Duv}
(-1)^{\frac{l(l-1)}{2}}D(\boldsymbol u,\boldsymbol v) = \det\big( a_{u_i+i+v_{l-j+1}+l-j} \big)_{i,j=1}^l>0,
\end{align}
where the last conclusion follows from the fact the determinant in question is a minor of the Toeplitz matrix \( (a_{i-j})_{i,j=1}^\infty \) corresponding to rows \( u_1+1+v_l+l,\ldots,u_l+l+v_l+l\) and columns \( 1,v_l-v_{l-1}+2,\ldots,v_l-v_1+l\). Finally, \eqref{sign Duv} together with \eqref{W Taylor} and \eqref{sign Wu} now shows that 
\[
(-1)^{\frac{l(l-1)}{2}} W_l(x,y)>0,  \quad x,y\in(0,r_{\mathrm r}). \qedhere
\]
\end{proof}

\end{document}
